\rok{Новембар 2024.}{D}

Први практични тест из Алгоритама и структура података

\zadatak{1}{Selection Sort}
Написати \textit{Selection Sort} алгоритам за сортирање целобројног низа дужине \textit{n}.

\textbf{Додатне напомене за решавање задатка:}
\begin{itemize}
    \item Улаз је несортиран низ целих бројева.
    \item Излаз је сортирани низ целих бројева.
    \item У шаблону се налази класа \texttt{RandomArrayGenerator} коју треба искористити за генерисање низа случајних бројева.
    \item Креирати класу \texttt{Selection} и у оквиру ње генерисати методу:
    \begin{Verbatim}[fontsize=\footnotesize, numbers=left, xleftmargin=10mm]
public static void Sort(long[] arr)
    \end{Verbatim}
\end{itemize}



\zadatak{2}{Минималан збир близанаца}
У повезаној листи са \textbf{парним бројем чворова} величине $n$, дефинишемо \textbf{близаначки пар} чворова на следећи начин:

За \textbf{i-ти чвор} (почевши од индекса 0) постоји одговарајући \textbf{близанац} на позицији $(n-1-i)$, ако важи $0 \leq i \leq (n/2)-1$.

На пример:
\begin{itemize}
    \item Када је $n=4$, чвор са индексом 0 има близанца у чвору са индексом 3, а чвор са индексом 1 има близанца у чвору са индексом 2. Ова два пара су једини близаначки парови за листу величине 4.
\end{itemize}

\textbf{Збир близанаца} дефинише се као збир вредности чвора и његовог близанца. Задатак је да пронађете \textbf{најмањи збир близаначких парова} у листи. Решење треба да пронађе минималан збир близанаца у $O(n)$ времену.

\textbf{Додатни захтеви за решавање задатка:}
\begin{itemize}
    \item Алгоритам писати у оквиру класе \texttt{LinkedList}.
    \item Захтеван потпис методе:
    \begin{Verbatim}[fontsize=\footnotesize, numbers=left, xleftmargin=10mm]
public int findMinTwinSum()
    \end{Verbatim}
\end{itemize}

\textbf{Примери:}
\begin{itemize}
    \item \textbf{Пример 1:} \\
          \textbf{Улаз:} $5 \leftrightarrow 4 \leftrightarrow 1 \leftrightarrow 1$ \\
          \textbf{Излаз:} 5
    \item \textbf{Пример 2:} \\
          \textbf{Улаз:} $4 \leftrightarrow 2 \leftrightarrow 2 \leftrightarrow 3$ \\
          \textbf{Излаз:} 4
\end{itemize}



\zadatak{3}{Верификација могућности пермутације излазног низа помоћу стека}
Креирати алгоритам који проверава да ли је могуће добити излазни низ из улазног низа користећи само стек. Улазни низ садржи јединствене бројеве, а излазни низ је потенцијална пермутација тих бројева коју желимо да добијемо.

\textbf{Додатни захтеви за решавање задатка:}
\begin{itemize}
    \item Задатак решавати помоћу структуре стека (класа \texttt{Stack}) која је дата у оквиру шаблона задатка.
    \item Захтеван потпис методе:
    \begin{Verbatim}[fontsize=\footnotesize, numbers=left, xleftmargin=10mm]
public static bool CanAchieveOutput(int[] input, int[] output)
    \end{Verbatim}
\end{itemize}

\textbf{Примери:}
\begin{itemize}
    \item \textbf{Улазни низ:} \texttt{[1, 2, 3, 4, 5]}, \textbf{Излазни низ:} \texttt{[4, 5, 3, 2, 1]}, \textbf{Испис:} \texttt{true}
    \item \textbf{Улазни низ:} \texttt{[1, 2, 3, 4, 5]}, \textbf{Излазни низ:} \texttt{[4, 3, 5, 1, 2]}, \textbf{Испис:} \texttt{false}
\end{itemize}

\sablon{Шаблон првог теста новембар 2024. група D}{2024/Novembar/GrupaD/AISP2024_Test1_GrupaD.linq}
